\scanpage{13}
\mspage{11}
se définir à l'aide d'un ``chemin'' dans $\deleted{\unsure{\operatorname{Hom}_\Sigma(}}\underline{\mathrm{Ch}}_{\Sigma,m'+n'}(X;x,z)$. Dans le cas général, on va en fait trouver une homotopie \deleted{\unsure{$k$}}
\[
h=h_{(n,m),(n',m')},
\]
``rendant commutatif'' le diagramme de foncteurs
\diagram{\begin{tikzcd}[ampersand replacement=\&,column sep=2.4em,scale=.85,transform shape,row sep=3.3em,cells={font=\scriptsize}]
A=\underline{\mathrm{Ch}}_{\Sigma,m}(X;x,y)\times\underline{\mathrm{Ch}}_{\Sigma,n}(X;y,z)
\arrow[r,"{*_{m,n}}"]\arrow[d,"{i_{m,m'}\times j_{\unsure{m|n},n'}}","{\deleted{\unsure{i_{m}}}}"']
\& \underline{\mathrm{Ch}}_{\Sigma,m+n}(X;x,z)\arrow[d,"{k_{m+n,m'+n'}}",""{name=R,inner sep=0pt}]\\
\underline{\mathrm{Ch}}_{\Sigma,m'}(X;x,y)\times\underline{\mathrm{Ch}}_{\Sigma,n'}(X;y,z)
\arrow[r,"{*_{m',n'}}",""{name=D,inner sep=0pt}]
\& \underset{\Vert}{\underline{\mathrm{Ch}}_{\Sigma,m'+n'}(X;x,z)}\arrow[from=R,to=D,bend right=17,"h"']\\
\& B
\end{tikzcd}}
où les $i_{\alpha,\beta},j_{\alpha,\beta},k_{\alpha,\beta}$ sont les foncteurs de transition
\[
\underline{\mathrm{Ch}}_{\Sigma,\alpha}(X;a,b)\,\unsure{\text{---}|\longrightarrow}\,\underline{\mathrm{Ch}}_{\Sigma,\beta}(X;a,b)
\]
$(\alpha\leq\beta)$ pour les \deleted{couples} $(x,y),(y,z)$ et $(x,z)$ (en guise de $(a,b)$) respectivement. L'homotopie $h$ sera \unsure{un|une} \deleted{suite de \unsure{flèches}} chemin dans la catégorie $\underline{\mathrm{Hom}}(A,B)$, \deleted{reliant} allant de
\[
k_{m+n,m'+n'}\circ *_{m,n}\longrightarrow *_{m',n'}\circ(i_{m,\unsure{m|m'}}\times j_{n,n'}).
\]
L'existence d'une telle homotopie signifie \deleted{\unsure{q}} ni plus ni moins que le diagramme carré ci-dessus est commutatif, en tant que diagramme dans la catégorie $(\mathrm{Cathot})$ des ``catégories à homotopie près'', dont les objets sont

% archive sheet number: 12
